\documentclass[11pt,a4paper]{amsart}

\usepackage[T1]{fontenc}
\usepackage[utf8]{inputenc}
\usepackage{amsmath,amssymb,amsthm,mathtools}
\usepackage{enumitem}
\usepackage{microtype}
\usepackage[margin=2.7cm]{geometry}
\usepackage{xcolor}
\definecolor{linkcol}{rgb}{0.10,0.20,0.50}
\usepackage[colorlinks=true,linkcolor=linkcol,citecolor=linkcol,urlcolor=linkcol]{hyperref}
\makeatletter
\g@addto@macro\UrlBreaks{\do\_\do\.}
\makeatother

\theoremstyle{plain}
\newtheorem{theorem}{Theorem}[section]
\newtheorem{proposition}[theorem]{Proposition}
\newtheorem{lemma}[theorem]{Lemma}
\newtheorem{corollary}[theorem]{Corollary}
\theoremstyle{definition}
\newtheorem{definition}[theorem]{Definition}
\newtheorem{example}[theorem]{Example}
\newtheorem{remark}[theorem]{Remark}

\newcommand{\CC}{\mathbb C}
\newcommand{\NN}{\mathbb N}
\newcommand{\QQ}{\mathbb Q}
\newcommand{\ZZ}{\mathbb Z}
\newcommand{\Ram}{\operatorname{Ram}}
\newcommand{\Irr}{\operatorname{Irr}}
\newcommand{\Slope}{\operatorname{Slope}}
\newcommand{\Spec}{\operatorname{Spec}}
\newcommand{\Loop}{\operatorname{Loop}}
\newcommand{\ord}{\operatorname{ord}}
\newcommand{\lcm}{\operatorname{lcm}}
\newcommand{\den}{\operatorname{den}}
\newcommand{\muR}{\boldsymbol\mu}
\emergencystretch=1.5em

\begin{document}

\title[Formal-contact covariance in Stokes geometry]{Formal-contact covariance, zero-loop rigidity\
and exact recognition for ramified nonabelian Hodge diagrams}

\author{M.~Haufschild}
\address{Independent researcher}
\email{haufschildmichael@gmail.com}

\date{3 August 2026}

\subjclass[2020]{Primary 34M40, 14D24; Secondary 14H60, 32G34, 57K10}
\keywords{irregular nonabelian Hodge theory, Stokes geometry, fission tree, Eggers--Wall tree,
contact complexity, Puiseux characteristic, ultrametric, Fourier transform}

\begin{abstract}
We identify a coordinate-covariance law governing loops, edges, Fourier rank drops and information
loss in ramified type-$A$
nonabelian Hodge diagrams.  For core matrix $B$ and ramification decoration $r_i$, put
\[
 K_{ij}=1+\frac{B_{ij}}{r_ir_j}\ (i\ne j),\qquad
 K_{ii}=1+\frac{B_{ii}-1}{r_i^2}.
\]
At source, this is the normalized survival integral of a cyclic Galois sheet-pair filtration:
irregularity is the coarea of the heights at which pairs of sheets cease to agree.
Universally, every tangent formal reparametrization $t\mapsto t u(t)$ with $u(0)=1$ preserves all
Laurent valuations, every pairwise contact order $\ord(f-g)$ and the complete threshold hierarchy.
For a negative Airy Fourier step, the reciprocal transformed coordinate is exactly
$t^E u(t)^{-1}=\phi(t^E)$ for such an automorphism $\phi$; hence the Fourier/contact statements are
instances of the same covariance law rather than parallel coincidences.
On the common-truncation quotient of the twisted fission tree, $K_{ij}$ is the complementary
Eggers--Wall contact coordinate of the meet of $i$ and $j$: raw exponent length is rescaled by the
reciprocal ramification index.  This reparametrization is strictly increasing on every branch.
Consequently the decorated diagram reconstructs the full labelled multi-circle branching
skeleton, not merely an ultrametric inequality.  Moreover
\[
 d_{ij}=2K_{ij}-K_{ii}-K_{jj}
\]
is its additive path metric; in particular it satisfies the four-point condition and has strict
negative type.  We obtain the new strict local obstruction $K_{ii}<K_{ij}$ for all $i\ne j$.

The diagonal critical value $K_{ii}=1-r_i^{-2}$ yields a rigid zero-loop stratum.  If a Stokes
circle has characteristic exponents $\beta_1/r>\cdots>\beta_m/r$ and cumulative gcd chain
$r=e_0>\cdots>e_m=1$, the sign of its loop is the sign of the first nonzero
$\beta_i-(e_{i-1}+e_i)$.  Thus zero loops are exactly the Airy towers
$\beta_i=e_{i-1}+e_i$, classified by strict divisor chains, equivalently ordered multiplicative
factorizations of $r$.  We compute leading-separated cores, exact elementary two-circle spectra,
and finite decorated fibres.  Pair-spectrum membership together with cophenetic consistency is
then necessary and sufficient for arbitrary finite elementary loopless decorated graphs.  The
proof constructs one simultaneous family of complex Puiseux polynomials and uses exactly
$n^2+n^3\le2n^3$ unit-cost recognition queries.
\end{abstract}

\maketitle

\section{Introduction}\label{sec:intro}

Wild nonabelian Hodge theory attaches hyperk\"ahler moduli spaces to irregular connections and
meromorphic Higgs bundles on curves \cite{BiquardBoalch2004}.  On the affine line, formal
singularity data determine a diagram whose symmetric integer matrix is a generalized Cartan matrix
for the corresponding nonabelian Hodge space \cite{BoalchYamakawa2020}.  Untwisted irregular types
give fission graphs.  Ramification changes the picture: loops and negative entries may occur, and
even a loopless nonnegative diagram need not be a fission graph.

Two recent developments make a sharp local analysis possible.  First, fission trees give the
discrete classification of admissible deformation classes, including arbitrary twisted type-$A$
formal structures \cite{BoalchDoucotRembado2025}.  Second, Dou\c cot wrote the loop and edge
multiplicities directly in terms of Puiseux exponents and proved that, after division by the
ramification products, all edge matrices satisfy ultrametric inequalities
\cite{Doucot2025Graphs}.  In the untwisted case those inequalities characterize fission graphs; in
the twisted case they are only necessary.

The organizing principle of this paper is that the loop and edge formulas are two evaluations of
one contact coordinate.  For the ramification decoration $r_i=\Ram(I_i)$, define
\begin{equation}\label{eq:intro-K}
 K_{ij}=1+\frac{B_{ij}}{r_ir_j}\quad(i\ne j),\qquad
 K_{ii}=1+\frac{B_{ii}-1}{r_i^2}.
\end{equation}
The fission tree is metrized by raw Puiseux exponent.  On a common truncation branch, let $R(u)$
be the cumulative ramification index at height $u$.  If $\eta$ lies above all relevant exponents,
the complementary contact coordinate is
\begin{equation}\label{eq:intro-contact}
 \kappa(v)=\eta-\int_{h(v)}^\eta\frac{du}{R(u)}.
\end{equation}
With the conventional Eggers--Wall orientation, the integral in
\eqref{eq:intro-contact} is contact complexity and $\kappa$ is its affine complement
\cite{GarciaBarrosoGonzalezPopescu2019}.  The shift is forced by measuring pole order downwards
from infinity rather than characteristic exponent upwards from a smooth reference branch.

The law above has a formal-coordinate source which is more general than the Stokes application.
Write $\CC((t))$ for Laurent series and $\ord$ for its valuation.

\begin{theorem}[Universal formal-contact covariance]\label{thm:intro-covariance}
Let $E\ge1$ and $u\in\CC[[t]]$ satisfy $u(0)=1$.  There is a tangent
$\CC$-algebra automorphism $\phi_u$ of $\CC((t))$, induced by
$t\mapsto t u(t)^{1/E}$, such that for all $f,g\in\CC((t))$
\[
 \ord(\phi_u f)=\ord(f),\qquad
 \ord(\phi_u f-\phi_u g)=\ord(f-g).
\]
Consequently simultaneous application of $\phi_u$ or $\phi_u^{-1}$ preserves every pairwise
contact matrix, every threshold cluster and the complete contact hierarchy of any indexed Laurent
family.

For a normalized Airy leading term $f=t^{-(R+E)}v(t)$ with $v(0)=1$, the local Fourier coordinate
$\widehat z=f/t^{-R}$ has order $-E$, while its reciprocal satisfies
\[
 \widehat\zeta=t^E v(t)^{-1}=\phi_{v^{-1}}(t^E).
\]
At operator level the genuine Weyl-algebra Fourier automorphism sends
$D^R-X^E$ to $(-X)^R-D^E$.  Iterating these exact operator and Laurent steps along an Airy tower
gives a strictly decreasing rank chain ending at one, while preserving contact tomography at every
step.
\end{theorem}

\begin{proof}
The binomial $E$th root of $u$ has constant term one.  Substitution by
$t u(t)^{1/E}$ has the inverse obtained from formal compositional inversion; localization extends
the substitution from $\CC[[t]]$ to $\CC((t))$.  Factoring a nonzero Laurent series into its leading
monomial and a unit proves valuation invariance, and applying this to $f-g$ proves contact
invariance.  The Fourier coordinate identities are direct Laurent multiplication and inversion.
The Weyl identity follows from $X\mapsto D$, $D\mapsto-X$.  Recursive application to the exact
characteristic tail produces the finite resolution.
\end{proof}

\begin{theorem}[Contact tomography; informal form]\label{thm:intro-tomography}
Replace the metrised fission tree by its common-truncation skeleton, gluing branches at their
fission exponents.  Then
\[
                         K_{ij}=\kappa(i\wedge j).
\]
The coordinate $\kappa$ is strictly increasing from every leaf to the root.  Hence $(B,r)$
reconstructs the labelled rooted branching skeleton after unary vertices and vertex types are
forgotten.  In addition,
\begin{equation}\label{eq:intro-energy}
 d_{ij}=\frac{2B_{ij}}{r_ir_j}-\frac{B_{ii}-1}{r_i^2}
        -\frac{B_{jj}-1}{r_j^2}
\end{equation}
is the additive path metric of this skeleton and is of strict negative type.
\end{theorem}

This theorem sharpens both parts of the known rescaled ultrametric statement.  Strictness along
leaf edges gives
\begin{equation}\label{eq:intro-strict-local}
 \frac{B_{ii}-1}{r_i^2}<\frac{B_{ij}}{r_ir_j}\qquad(i\ne j),
\end{equation}
while strictness at nested branch vertices shows that averaging over Galois sheets never collapses
the multi-circle hierarchy.  Thus the diagram loses characteristic and sibling decorations, but
not the underlying branching topology.  The path metric \eqref{eq:intro-energy} packages every
loop and edge into one cut-metric decomposition and supplies four-point and spectral constraints.

Against this universal background, the simplest unresolved local question is:
\begin{quote}
Which ramified Stokes circles have zero diagonal multiplicity, and which loopless decorated graphs
do they realize?
\end{quote}

The answer is unexpectedly rigid.  Let $I=\langle q\rangle$ have ramification $r>1$.  Retain only
the exponents that decrease the cumulative gcd, write them as
\[
  \frac{\beta_1}{r}>\cdots>\frac{\beta_m}{r}>0,
  \qquad r=e_0>e_1>\cdots>e_m=1,
  \qquad e_i=\gcd(e_{i-1},\beta_i).
\]
The diagonal entry is
\begin{equation}\label{eq:intro-loop}
  B_{I,I}=\sum_{i=1}^m\beta_i(e_{i-1}-e_i)-r^2+1.
\end{equation}
Our first theorem is a lexicographic sign law.

\begin{theorem}[Airy-tower sign law]\label{thm:intro-sign}
Let $j$ be the first index such that $\beta_j\ne e_{j-1}+e_j$.  Then
\[
 \operatorname{sgn}(B_{I,I})=
 \operatorname{sgn}\bigl(\beta_j-e_{j-1}-e_j\bigr).
\]
If there is no such index, then $B_{I,I}=0$.  In particular
\begin{equation}\label{eq:intro-zero}
 B_{I,I}=0\quad\Longleftrightarrow\quad
 \beta_i=e_{i-1}+e_i\quad(1\le i\le m).
\end{equation}
\end{theorem}

The proof is not an estimate near equality.  The gcd constraint creates a forbidden interval:
below $R+E$ the exponent is at most $R-E$, while above it the exponent is at least $R+2E$.
A subcanonical first term leaves a deficit that no decreasing tail can repair; a supercanonical
first term already exceeds the target $R^2-1$; and the canonical term
$(R+E)(R-E)=R^2-E^2$ reduces the problem exactly from $R$ to $E$.

We call the zero-loop data in \eqref{eq:intro-zero} \emph{Airy towers}.  Choosing an Airy tower is
equivalent to choosing the divisor chain $e_0>\cdots>e_m=1$, hence to choosing the ordered
factorization
\[
  r=a_1\cdots a_m,\qquad a_i=e_{i-1}/e_i\ge2.
\]
This arithmetic classification has three geometric effects.

First, it identifies a large Fourier-resolvable family.  The equality
$\beta_i=e_{i-1}+e_i$ is exactly the negative congruence
$e_{i-1}\equiv-e_i\pmod{\beta_i}$ in Hiroe's dual-Puiseux criterion
\cite{Hiroe2017}; every rank drop preserves the zero loop.

Second, it exposes information forgotten by the diagram.  For a leading-separated family, put
$a_i=r_i/e_{i,1}$.  We prove
\begin{equation}\label{eq:intro-separated}
  B_{I_i,I_j}=\frac{r_ir_j}{\min(a_i,a_j)}.
\end{equation}
The matrix depends only on the first factors $a_i$, while the remaining ordered factorization is
retained by the fission tree.  Thus one graph may carry a large, explicitly counted family of
admissible deformation classes.

Third, even the simplest ramified vertices impose congruence gaps on edges.  A simple Airy circle
of ramification $r$ has the single mandatory level $(r+1)/r$.  For two such circles the exact
spectrum is
\begin{equation}\label{eq:intro-spectrum}
 \Spec(r,s)=
 \begin{cases}
   \{\max(r,s)\}\cup rs\NN_{>0},&r\ne s,\\
   \{0,1,\ldots,r\}\cup r^2\NN_{>0},&r=s.
 \end{cases}
\end{equation}
No three-vertex test can see these two-vertex gaps.  Combining \eqref{eq:intro-spectrum} with the
known rescaled ultrametric inequality yields our global recognition theorem.

\begin{theorem}[Elementary recognition; informal form]\label{thm:intro-recognition}
Let $B$ be a symmetric nonnegative integer matrix with zero diagonal on a finite labelled set $N$,
and let $r_i\ge1$.  The decorated graph $(B,r)$ is the core diagram of elementary loopless Stokes
circles if and only if
\[
 B_{ij}\in\Spec(r_i,r_j)\quad(i\ne j)
\]
and the rational weights $w_{ij}=B_{ij}/(r_ir_j)$ have no triangle with a unique largest side.
The proof constructs one simultaneous family of finite complex Puiseux polynomials.  Its exhaustive
unit-cost schedule has exactly $|N|^2+|N|^3\le2|N|^3$ pair and triangle queries for $N\ne\varnothing$.
\end{theorem}

Here ``elementary'' means unramified when $r_i=1$ and simple Airy when $r_i>1$.  Sufficiency is the
new point.  A finite ultrametric dendrogram first realizes integral weights by common unramified
terms.  It then processes the smallest remaining ramification: mandatory Airy coefficients split
the leaves into root-of-unity orbits, with all larger ramifications occupying the unique empty
child allowed by the fission-tree axioms.  Lower fractional coefficients realize the residual
values $0,\ldots,r-1$.  This exhausts the exact pair spectra.

The result is deliberately local to formal Stokes geometry.  It does not classify all twisted
nonabelian Hodge graphs, nor assert that equal core diagrams give isomorphic wild character
varieties.  Recent work shows precisely why such an inference would be unsafe: fission data refine
the diagram and can distinguish admissible deformation classes invisible in the Cartan shadow
\cite{Doucot2026Basic}.  Our theorems instead identify and reconstruct one natural ramified stratum
exactly.

\subsection*{Organization}
Section~\ref{sec:background} fixes conventions and recalls the source formulas.
Section~\ref{sec:contact} proves the complementary contact law, strict tomography and the energy-tree
theorem.  Section~\ref{sec:sign} proves the sign law.  Section~\ref{sec:classification} gives the
divisor-chain, enumerative and Fourier consequences.  Sections~\ref{sec:separated}--\ref{sec:finite}
treat separated cores, pair spectra and finite fibres.  Section~\ref{sec:recognition} proves exact
elementary recognition.  Section~\ref{sec:examples} records graph twins and obstructions.
Section~\ref{sec:scope} records scope limits.

\section{Stokes circles, levels and diagrams}\label{sec:background}

We use the conventions of \cite[\S2]{Doucot2025Graphs}.  Let $z$ be the affine coordinate near
infinity.  A Stokes circle is a Galois orbit $I=\langle q\rangle$ of a Puiseux polynomial
\[
  q=\sum_{j=1}^{s}c_jz^{j/r},\qquad c_s\ne0,
\]
where $r$ is minimal.  Then $\Ram(I)=r$, $\Irr(I)=s$, and $\Slope(I)=s/r$.  Exponents whose
successive adjunction strictly lowers the cumulative gcd are the levels.  Deleting all other
exponents preserves the loop multiplicity.

For a ramified circle, list the level numerators in decreasing order as
\begin{equation}\label{eq:char-data}
 \beta_1>\cdots>\beta_m>0,
 \qquad e_0=r,
 \qquad e_i=\gcd(r,\beta_1,\ldots,\beta_i).
\end{equation}
Since each entry is a level,
\begin{equation}\label{eq:gcd-chain}
 r=e_0>e_1>\cdots>e_m=1,
 \qquad e_i=\gcd(e_{i-1},\beta_i).
\end{equation}
The last equality $e_m=1$ is equivalent to minimality of $r$.

\begin{definition}\label{def:weighted-loop}
For characteristic data \eqref{eq:char-data}--\eqref{eq:gcd-chain}, define
\[
 S(\beta,e)=\sum_{i=1}^m\beta_i(e_{i-1}-e_i),
 \qquad
 \Loop(\beta,e)=S(\beta,e)-r^2+1.
\]
\end{definition}

The loop formula of \cite[Lemma~2.1]{Doucot2025Graphs} says
\begin{equation}\label{eq:source-loop}
 B_{I,I}=\Loop(\beta,e).
\end{equation}
The sum $S$ is also the Komatsu--Malgrange irregularity of the endomorphism connection, in the
dual-Puiseux form developed by Hiroe \cite{Hiroe2017}.  Formula \eqref{eq:source-loop} is the only
input to the arithmetic argument below.

The following source lemma makes the passage from Galois sheets to $S$ explicit.  Its role is to
close the source chain, not to assert novelty over the conductor calculation implicit in the
published loop formula.

\begin{proposition}[Cyclic stabilizer coarea law]\label{prop:cyclic-source}
Let
\[
 q_{\rm char}=\sum_{i=1}^{m}z^{\beta_i/r},
 \qquad \zeta_r=\exp(2\pi\sqrt{-1}/r),
 \qquad \sigma_a(q_{\rm char})=\sum_i\zeta_r^{a\beta_i}z^{\beta_i/r}
\]
for $a\in\ZZ/r\ZZ$.  If $P_i=(\beta_1,\ldots,\beta_i)$ and
\[
 H_i=\{a:\sigma_a\text{ fixes every monomial indexed by }P_i\},
 \qquad H_0=\ZZ/r\ZZ,
\]
then $|H_i|=e_i$.  The sheets whose first difference occurs at
$\beta_i/r$ form $H_{i-1}\setminus H_i$; hence
\[
 \sum_{a,b\in\ZZ/r\ZZ}\Slope\bigl(\sigma_a(q_{\rm char})-
 \sigma_b(q_{\rm char})\bigr)
 =\sum_i\beta_i(e_{i-1}-e_i)=S(\beta,e).
\]
Moreover $e_m=1$ makes the translate map injective, so its Galois orbit has exactly $r$ sheets, and
the loop reconstructed from these ordered sheet differences is exactly
$S(\beta,e)-r^2+1$.
\end{proposition}

\begin{proof}
Primitivity of $\zeta_r$ gives $\zeta_r^{aj}=1$ precisely when $r\mid aj$.  If $o(a)$ is the
additive order of $a$, fixing every exponent in $P_i$ is therefore equivalent to
$o(a)\mid\gcd(r,P_i)=e_i$.  Thus $H_i$ is the kernel of multiplication by $e_i$ on the cyclic group
of order $r$, and its order is $e_i$.  The first-difference assertion follows by subtracting two
successive nested stabilizers.  For one base sheet its difference of level $\beta_i/r$ occurs
$e_{i-1}-e_i$ times; translation gives $r$ copies among ordered pairs, cancelling the denominator
$r$.  Finally $e_m=1$ makes the stabilizer trivial, so orbit--stabilizer and the stated loop
normalization give the last assertions.
\end{proof}

For two distinct circles $I,J$, their common part is the maximal common Galois-orbit truncation,
and the fission exponent $f_{I,J}$ is the larger slope of the two remaining parts.  The exact edge
formula \cite[Lemma~2.2]{Doucot2025Graphs} has one contribution from each level of the common part,
one contribution from $f_{I,J}$, and the normalization $-\Ram(I)\Ram(J)$.  We will use two direct
specializations.  If the common part is zero and the larger slope is $s/r$, then
\begin{equation}\label{eq:zero-common-edge}
 B_{I,J}=r'(s-r),
\end{equation}
where $r'$ is the ramification of the other circle.  If both circles are unramified, then
\begin{equation}\label{eq:unramified-edge}
 B_{I,J}=\deg(q-q')-1.
\end{equation}

The diagonal Proposition~\ref{prop:cyclic-source} and the edge formula are instances of one finite
coarea law.  We record its two-orbit form because it is the source of the contact coordinate used
throughout the paper.

\begin{theorem}[Galois sheet-pair coarea law]\label{thm:sheet-pair-coarea}
Write the nonzero common exponents in decreasing order as reduced fractions
$k_i=n_i/d_i$, put
\[
 D_{-1}=1,\qquad D_i=\mathop{\rm lcm}(d_0,\ldots,d_i),
 \qquad m_i=\frac r{d_i}n_i,
\]
and suppose the fission exponent is $s_d/r$, with this side having the larger different-part
slope.  If $E_i$ is the set of ordered sheet pairs in
$(\ZZ/r\ZZ)\times(\ZZ/r'\ZZ)$ which agree through the first $i$ common exponents, then
\[
 |E_i|=\frac{rr'}{D_i},
 \qquad
 |E_{i-1}\setminus E_i|
 =r\left(\frac{r'}{D_{i-1}}-\frac{r'}{D_i}\right).
\]
Consequently
\begin{equation}\label{eq:pair-coarea-source}
 \Irr\!\operatorname{Hom}(I,J)
 =\sum_i m_i\left(\frac{r'}{D_{i-1}}-\frac{r'}{D_i}\right)
   +s_d\frac{r'}{D_t},
 \qquad
 B_{I,J}=\Irr\!\operatorname{Hom}(I,J)-rr'.
\end{equation}
Together with Proposition~\ref{prop:cyclic-source}, every diagonal and off-diagonal core entry is
obtained from the same loss-of-sheet-equality coarea.
\end{theorem}

\begin{proof}
At the reduced exponent $n/d$, the phases are
$\exp(2\pi\sqrt{-1}an/d)$ and $\exp(2\pi\sqrt{-1}bn/d)$.  Coprimality of $n,d$ makes them equal
exactly when $a\equiv b\pmod d$.  Agreement through a prefix is therefore congruence modulo the
least common denominator $D_i$.  The difference homomorphism
\[
 (\ZZ/r\ZZ)\times(\ZZ/r'\ZZ)\longrightarrow\ZZ/D_i\ZZ,
 \qquad(a,b)\longmapsto a-b,
\]
is surjective, so its kernel has order $rr'/D_i$.  Subtracting consecutive kernels gives the layer
size.  Multiplication by the common slope $m_i/r$ cancels $r$; the surviving kernel separates at
the fission slope $s_d/r$.  Summing proves the irregularity identity, and the standard
$-rr'$ normalization gives the edge identity.
\end{proof}

For a finite irregular class $\Theta=\sum_{i\in N}n_iI_i$, these numbers form the symmetric core
matrix $B=(B_{ij})$.  The rescaled off-diagonal entries are
\begin{equation}\label{eq:rescaled}
 \widetilde B_{ij}=\frac{B_{ij}}{r_ir_j}\qquad(i\ne j).
\end{equation}
On the diagonal the compatible normalization is
\begin{equation}\label{eq:rescaled-diagonal}
 \widetilde B_{ii}=\frac{B_{ii}-1}{r_i^2}.
\end{equation}
Every decorated nonabelian Hodge diagram satisfies the ultrametric implication
\begin{equation}\label{eq:known-ultra}
 \widetilde B_{ij}\le\widetilde B_{ik}\le\widetilde B_{jk}
 \quad\Longrightarrow\quad
 \widetilde B_{ik}=\widetilde B_{jk}
\end{equation}
for pairwise distinct $i,j,k$, together with
\begin{equation}\label{eq:known-diagonal}
 \widetilde B_{ii}\le \widetilde B_{ij}\qquad(i\ne j),
\end{equation}
by \cite[Theorem~4.2]{Doucot2025Graphs}.  Section~\ref{sec:contact} proves that
\eqref{eq:known-diagonal} is always strict and identifies exactly which part of the fission tree is
retained by the full matrix of rescaled entries.

\section{Complementary contact complexity and strict tomography}\label{sec:contact}

We first rewrite the source formulas in their invariant tree form.  This is a bridge, rather than a
novelty claim: contact complexity is standard on Eggers--Wall trees, and the finite sum below is
equivalent to the rescaled formulas of \cite[Lemma~4.3]{Doucot2025Graphs}.  The new content begins
with the strict comparisons in Lemmas~\ref{lem:strict-diagonal} and
\ref{lem:strict-nested} and their reconstruction consequences.

Let $I_i,I_j$ be distinct circles.  Write the levels of their common part as
\[
 \lambda_0>\lambda_1>\cdots>\lambda_t>f_{ij},
\]
where $f_{ij}=f_{I_i,I_j}$ and the level list may be empty.  Put
\[
 D_a=\lcm(\den(\lambda_0),\ldots,\den(\lambda_a)),\qquad D_{-1}=1.
\]
Here $D_a$ is the ramification index of the corresponding common truncation.  Define
\begin{equation}\label{eq:contact-pair-sum}
 K_{ij}=\lambda_0(1-D_0^{-1})+
 \sum_{a=1}^t\lambda_a(D_{a-1}^{-1}-D_a^{-1})+f_{ij}D_t^{-1};
\end{equation}
when the level list is empty, the right side means $f_{ij}$.  For $i=j$, use all levels of $I_i$,
put $f_{ii}=0$, and take the terminal denominator to be $r_i$.

Choose an integer $\eta$ above all the levels under consideration.  On the common branch let
$R(u)=1$ for $u\in(\lambda_0,\eta)$ and $R(u)=D_a$ between consecutive levels, with
$R(u)=D_t$ on $(f_{ij},\lambda_t)$.  Summation by parts gives
\begin{equation}\label{eq:contact-integral}
\begin{aligned}
 \eta-\int_{f_{ij}}^\eta\frac{du}{R(u)}
 &=\lambda_0-\frac{\lambda_0-\lambda_1}{D_0}-\cdots
   -\frac{\lambda_t-f_{ij}}{D_t}\\
 &=K_{ij}.
\end{aligned}
\end{equation}
For the diagonal the same calculation gives
\begin{equation}\label{eq:contact-diagonal-integral}
 K_{ii}=\eta-\int_0^\eta\frac{du}{R_i(u)}.
\end{equation}
The value is independent of the sufficiently large choice of $\eta$.  With the usual
Eggers--Wall exponent directed away from its smooth reference branch, the integral is precisely
the contact-complexity coordinate $c=\int de/R$; our $K=\eta-c$ is its affine complement because
pole height is oppositely oriented \cite[Definition~3.15 and Theorem~3.25]
{GarciaBarrosoGonzalezPopescu2019}.

The rescaled source formulas now become one identity.

\begin{proposition}[Complementary contact law]\label{prop:contact-law}
For every finite irregular class at infinity,
\begin{equation}\label{eq:contact-B-law}
 K_{ij}=1+\widetilde B_{ij}
 =\begin{cases}
  1+B_{ij}/(r_ir_j),&i\ne j,\\
  1+(B_{ii}-1)/r_i^2,&i=j.
 \end{cases}
\end{equation}
Equivalently, $K_{ij}$ is the complementary contact coordinate at the meet of the $i$th and
$j$th branches, with a leaf viewed as its own meet.
\end{proposition}

\begin{proof}
Equations \eqref{eq:contact-pair-sum} and \eqref{eq:contact-integral} are exactly the edge formula
of \cite[Lemma~4.3]{Doucot2025Graphs} after adding $1$ to $\widetilde B_{ij}$.  The diagonal
formula \eqref{eq:contact-diagonal-integral} is its loop counterpart.  The common-truncation
interpretation follows because the ramification index at a point of the fission tree is the least
common multiple of the denominators of its mandatory ancestors
\cite[Remark~3.19]{BoalchDoucotRembado2025}.
\end{proof}

The positivity of reciprocal ramification turns the weak inequalities recalled above into strict
ones whenever the corresponding fission events are distinct.

\begin{lemma}[Strict diagonal separation]\label{lem:strict-diagonal}
For distinct Stokes circles $I_i,I_j$,
\begin{equation}\label{eq:strict-diagonal}
 K_{ii}<K_{ij},\qquad K_{jj}<K_{ij}.
\end{equation}
Equivalently,
\begin{equation}\label{eq:strict-integral-obstruction}
 r_iB_{ij}>r_j(B_{ii}-1),\qquad r_jB_{ij}>r_i(B_{jj}-1).
\end{equation}
\end{lemma}

\begin{proof}
Distinct circles have positive fission exponent $f_{ij}>0$.  Above $f_{ij}$, the index profile of
$I_i$ agrees with the common profile.  Subtracting \eqref{eq:contact-diagonal-integral} from
\eqref{eq:contact-integral} therefore gives
\[
 K_{ij}-K_{ii}=\int_0^{f_{ij}}\frac{du}{R_i(u)}>0.
\]
The same argument with $j$ proves the second inequality.  Clearing positive denominators gives
\eqref{eq:strict-integral-obstruction}.
\end{proof}

\begin{lemma}[Strict nested separation]\label{lem:strict-nested}
Let $i,j,k$ be distinct leaves.  If the common-truncation meet of $i,j$ lies strictly below the
common meet of $i,k$ and $j,k$, then
\begin{equation}\label{eq:strict-rooted-triple}
 K_{ij}<K_{ik}=K_{jk}.
\end{equation}
If the three pairwise meets coincide, the three off-diagonal $K$-values coincide.
\end{lemma}

\begin{proof}
Pairs separated into different children of one branch vertex have the same common levels and the
same fission exponent, so their values in \eqref{eq:contact-pair-sum} agree.

For strictness, cancel the common prefix ending with denominator $D_t$.  Write $f$ for the higher
fission exponent.  The lower meet adds possible levels $\lambda_{t+1},\ldots,\lambda_u\le f$ and
ends at a strictly smaller exponent $g<f$.  Its remaining contribution is
\[
 \sum_{a=t+1}^u\lambda_a(D_{a-1}^{-1}-D_a^{-1})+gD_u^{-1}.
\]
The nonnegative coefficients sum to $D_t^{-1}$.  Every sampled height is at most $f$, and the
positive final coefficient $D_u^{-1}$ samples $g<f$.  This strict convex combination is smaller
than the higher contribution $fD_t^{-1}$.
\end{proof}

Let $\mathcal S(\Theta)$ denote the labelled rooted \emph{branching skeleton}: retain the hierarchy
of the fission tree, suppress unary vertices, and forget mandatory, inconsequential and empty
vertex types.  If $v$ is an internal vertex, choose leaves in two different children and set
$\kappa(v)=K_{ij}$; put $\kappa(i)=K_{ii}$ at a leaf.  Lemma~\ref{lem:strict-nested} makes the
first definition independent of the chosen leaves.

\begin{theorem}[Strict contact tomography]\label{thm:contact-tomography}
The coordinate $\kappa$ is strictly increasing from each leaf towards the root and
\[
 K_{ij}=\kappa(i\wedge j).
\]
The decorated diagram $(B,r)$ reconstructs $\mathcal S(\Theta)$ uniquely as a labelled rooted
tree.  Consequently every fission tree in a fixed decorated fibre has the same branching
skeleton.
\end{theorem}

\begin{proof}
Strictness on leaf edges is Lemma~\ref{lem:strict-diagonal}; strictness between nested branch
vertices is Lemma~\ref{lem:strict-nested}.  For reconstruction, list the distinct off-diagonal
values of $K$ and, for each such value $a$, set
\[
 i\sim_a j\quad\Longleftrightarrow\quad K_{ij}\le a,
\]
with reflexivity imposed on the diagonal.  Lemma~\ref{lem:strict-nested} makes $\sim_a$ an
equivalence relation, and its classes form a nested hierarchy as $a$ increases.  Attach each class
to the least strictly larger class containing it and attach the labelled singleton leaves at the
bottom.  Every branching cluster of $\mathcal S(\Theta)$ occurs because $\kappa$ is strict along
root paths, and every reconstructed class is the descendant set of such a branch vertex.  The
resulting isomorphism is label preserving.  The final claim follows because $K$ is determined by
$(B,r)$.
\end{proof}

The theorem identifies exactly why the fission tree remains finer than the diagram
\cite{Doucot2026Basic}: rational raw heights, unary vertices, ramification profiles and sibling
types can vary over a fixed skeleton.  The multi-leaf branching hierarchy cannot.

\begin{theorem}[Stokes energy-tree law]\label{thm:energy-tree}
Set $d_{ii}=0$ and, for $i\ne j$,
\begin{equation}\label{eq:energy-metric}
 d_{ij}=2K_{ij}-K_{ii}-K_{jj}
 =\frac{2B_{ij}}{r_ir_j}-\frac{B_{ii}-1}{r_i^2}-\frac{B_{jj}-1}{r_j^2}.
\end{equation}
Then $d$ is the additive path metric on $\mathcal S(\Theta)$ with contact-coordinate edge lengths.
In particular:
\begin{enumerate}[label=(\roman*)]
\item among $d_{ij}+d_{kl}$, $d_{ik}+d_{jl}$ and $d_{il}+d_{jk}$, the two largest are equal;
\item $d$ has strict negative type;
\item there is an explicit Hilbert embedding whose squared distance is $d$.
\end{enumerate}
\end{theorem}

\begin{proof}
Give the edge from a child $v$ to its parent $p(v)$ length
$\ell_v=\kappa(p(v))-\kappa(v)>0$.  The path from $i$ to $j$ rises from $K_{ii}$ to $K_{ij}$ and
descends to $K_{jj}$, proving \eqref{eq:energy-metric}.  The four-point condition is the standard
edge cancellation for a tree metric.

For the remaining claims, cut the leaf set at every edge $e$, choose one side $A_e$, and write
$\ell_e>0$ for its length.  Then
\[
 d_{ij}=\sum_e\ell_e\left|\mathbf 1_{A_e}(i)-\mathbf 1_{A_e}(j)\right|.
\]
If $x\ne0$ and $\sum_i x_i=0$, then
\[
 \sum_{i,j}x_ix_jd_{ij}=-2\sum_e\ell_e\left(\sum_{i\in A_e}x_i\right)^2<0.
\]
Equality would force the sum on every positive pendant cut $\{i\}$ to vanish.  Finally,
$\Phi(i)_e=\sqrt{\ell_e}\,\mathbf1_{A_e}(i)$ satisfies
$\|\Phi(i)-\Phi(j)\|^2=d_{ij}$.
\end{proof}

\section{The first-deviation sign law}\label{sec:sign}

The proof rests on a discrete interval excluded by the gcd condition.

\begin{lemma}[Gcd gap]\label{lem:gcd-gap}
Let $R>E\ge1$, assume $E\mid R$, and let $\gcd(R,\beta)=E$.
Then:
\begin{enumerate}[label=(\roman*)]
\item $\beta=R+E$ satisfies the gcd condition;
\item if $\beta<R+E$, then $\beta\le R-E$;
\item if $\beta>R+E$, then $\beta\ge R+2E$.
\end{enumerate}
\end{lemma}

\begin{proof}
Write $R=aE$ and $\beta=cE$.  Then $a\ge2$ and $\gcd(a,c)=1$.  The value $c=a+1$ is coprime to
$a$.  If $c<a+1$, then $c\ne a$, hence $c\le a-1$.  If $c>a+1$, integrality gives $c\ge a+2$.
Multiplying by $E$ proves the result.
\end{proof}

We isolate the tail estimate used in the negative branch.

\begin{lemma}[First-exponent bound]\label{lem:tail-bound}
Let a nonempty characteristic datum start at ramification $R$ with first exponent $\beta$.
Then
\[
 S(\beta,e)\le \beta(R-1).
\]
If the first gcd drop is to an intermediate $E>1$, the inequality is strict.
\end{lemma}

\begin{proof}
The drops telescope:
\[
 \sum_i(e_{i-1}-e_i)=R-1.
\]
Every exponent is at most $\beta$, giving the weak inequality.  If a tail exists, all its
exponents are strictly less than $\beta$, and its total drop is $E-1>0$, giving strictness.
\end{proof}

\begin{lemma}[First-deviation trichotomy]\label{lem:first-deviation}
Let the first step of characteristic data be
\[
 R\longrightarrow E,\qquad \beta=\beta_1.
\]
Then:
\begin{enumerate}[label=(\roman*)]
\item if $\beta<R+E$, then $\Loop(\beta,e)<0$;
\item if $\beta>R+E$, then $\Loop(\beta,e)>0$;
\item if $\beta=R+E$, then $\Loop(\beta,e)$ equals the loop functional of the tail beginning at
$E$.
\end{enumerate}
\end{lemma}

\begin{proof}
Suppose first that $\beta<R+E$.  Lemma~\ref{lem:gcd-gap} gives $\beta\le R-E\le R-1$.
If there is no tail, then
\[
 S=\beta(R-1)\le(R-1)^2<R^2-1.
\]
If a tail exists, Lemma~\ref{lem:tail-bound} gives the strict inequality
$S<\beta(R-1)\le(R-1)^2<R^2-1$.  Thus $\Loop<0$.

Suppose next that $\beta>R+E$.  Lemma~\ref{lem:gcd-gap} gives $\beta\ge R+2E$.  Since $E$ is a
proper divisor of $R$, one has $R\ge2E$.  The first summand alone satisfies
\[
 \beta(R-E)\ge(R+2E)(R-E)
 =R^2+E(R-2E)\ge R^2>R^2-1.
\]
All other summands are nonnegative, so $\Loop>0$.

Finally, if $\beta=R+E$, then
\[
 (R+E)(R-E)=R^2-E^2.
\]
Substitution into Definition~\ref{def:weighted-loop} cancels $R^2-E^2$ and leaves precisely the
tail functional normalized by $E^2-1$.
\end{proof}

\begin{theorem}[Airy-tower sign law]\label{thm:sign-law}
For characteristic data \eqref{eq:char-data}--\eqref{eq:gcd-chain}, let $j$ be the first index such
that $\beta_j\ne e_{j-1}+e_j$.  If $j$ exists, then
\[
 \operatorname{sgn}B_{I,I}
 =\operatorname{sgn}\bigl(\beta_j-e_{j-1}-e_j\bigr).
\]
If $j$ does not exist, then $B_{I,I}=0$.  Hence
\[
 B_{I,I}=0\quad\Longleftrightarrow\quad
 \beta_i=e_{i-1}+e_i\quad\text{for all }i.
\]
\end{theorem}

\begin{proof}
Apply Lemma~\ref{lem:first-deviation}(iii) repeatedly while the characteristic value is canonical.
At the first noncanonical entry, parts (i) and (ii) give the strict sign.  If every entry is
canonical, the recursion ends at ramification $1$ with empty weighted sum, whose normalized loop is
$-1+1=0$.
\end{proof}

\begin{definition}
A ramified characteristic satisfying $\beta_i=e_{i-1}+e_i$ at every level is an
\emph{Airy tower}.
\end{definition}

\begin{remark}
Inconsequential exponents do not create additional zero-loop cases: they do not lower the
cumulative gcd and therefore do not occur in \eqref{eq:source-loop}.  An Airy tower classifies the
level data, not every coefficient in a representative Puiseux polynomial.
\end{remark}

\section{Divisor chains, enumeration and Fourier reduction}\label{sec:classification}

The zero-loop equality removes the exponents as independent variables.

\begin{theorem}[Arithmetic classification]\label{thm:factor-classification}
For $r>1$, the following sets are naturally bijective:
\begin{enumerate}[label=(\roman*)]
\item Airy towers of ramification $r$;
\item strict divisor chains
\[
 r=e_0>e_1>\cdots>e_m=1,\qquad e_i\mid e_{i-1};
\]
\item ordered multiplicative factorizations
\[
 r=a_1\cdots a_m,\qquad a_i\ge2.
\]
\end{enumerate}
\end{theorem}

\begin{proof}
Given a divisor chain, set $\beta_i=e_{i-1}+e_i$.  Then
\[
 \beta_i-\beta_{i+1}=e_{i-1}-e_{i+1}>0
\]
and, since $e_i\mid e_{i-1}$,
\[
 \gcd(e_{i-1},\beta_i)=\gcd(e_{i-1},e_i)=e_i.
\]
Thus the prescribed values form characteristic data, and
Theorem~\ref{thm:sign-law} makes them an Airy tower.  Conversely, that theorem forces these values
from any zero-loop characteristic, so the first correspondence is injective and surjective.

For the second, put $a_i=e_{i-1}/e_i$.  The factors are at least two because the chain is strict,
and they telescope to $r$.  Conversely, from $r=a_1\cdots a_m$ define
$e_i=a_{i+1}\cdots a_m$ (with $e_m=1$).  These operations are inverse.
\end{proof}

Let $A(r)$ denote the number of Airy towers of ramification $r$, and put $A(1)=1$ for the empty
tail.

\begin{corollary}[Two exact counts]\label{cor:counts}
\begin{enumerate}[label=(\roman*)]
\item If $r=p^k$ with $p$ prime and $k\ge1$, then $A(r)=2^{k-1}$.
\item If $r$ is the product of $k$ distinct primes, then
\[
 A(r)=\sum_{j=1}^k j!\,S(k,j),
\]
the $k$th ordered Bell number, where $S(k,j)$ is a Stirling number of the second kind.
\end{enumerate}
\end{corollary}

\begin{proof}
An ordered factorization of $p^k$ is a composition of $k$; its cuts are an arbitrary subset of the
$k-1$ interior gaps.  In the squarefree case, each factor is a block of prime factors, and the
ordered list of factors is exactly an ordered set partition of the $k$ primes.  An ordered
partition into $j$ blocks is counted by $j!S(k,j)$.
\end{proof}

The same equality controls the local Fourier transform at operator, Laurent-coordinate and contact
levels.

\begin{theorem}[Negative Fourier branch]\label{thm:fourier}
Every Airy tower satisfies
\begin{equation}\label{eq:negative-congruence}
 e_{i-1}\equiv-e_i\pmod{\beta_i}
\end{equation}
at every level.  At a step $R\to E$, the Weyl Fourier automorphism sends
$D^R-X^E$ to $(-X)^R-D^E$, and the reciprocal Laurent Fourier coordinate is tangent-conjugate to
$t^E$.  Recursion on the characteristic tail gives a finite exact negative-Fourier chain ending at
rank one.  Every step preserves all Laurent contact orders and the full threshold hierarchy, and
the diagonal multiplicity remains zero at each intermediate rank.

Conversely, if a characteristic satisfies \eqref{eq:negative-congruence} and
$\beta_i>e_{i-1}$ at every level, then it is an Airy tower.
\end{theorem}

\begin{proof}
For an Airy tower, \eqref{eq:negative-congruence} is the identity
$\beta_i=e_{i-1}+e_i$.  The exact Weyl, Laurent and contact assertions are the corresponding clauses
of Theorem~\ref{thm:intro-covariance}.  Applying the same step to each successive characteristic
tail gives the recursive descent.  Loop preservation is the canonical clause of
Lemma~\ref{lem:first-deviation}.

For the converse, the congruence says $\beta_i\mid e_{i-1}+e_i$.  Since
$e_i<e_{i-1}<\beta_i$, one has
\[
 \beta_i\le e_{i-1}+e_i<2\beta_i.
\]
The only positive multiple of $\beta_i$ in this interval is $\beta_i$ itself.  Hence
$\beta_i=e_{i-1}+e_i$ at every level.
\end{proof}

For comparison, Hiroe relates the same negative congruences to dual-Puiseux Fourier reduction and
iterated cabling \cite{Hiroe2017}.  No knot-theoretic statement is used in any theorem here.
Fourier invariance of the global nonabelian Hodge diagram is known in much greater generality, and
its role in a recent classification of type-$A$ Argyres--Douglas dualities is made explicit in
\cite{Doucot2026Fourier}.  The assertion above is different: it identifies the actual tangent
Laurent automorphism preserving every contact threshold and couples it to an exact Weyl-operator
resolution on the zero-loop Airy locus.

Topological local Laplace transforms of Stokes-filtered systems and explicit Fourier algorithms
for analytic Stokes data are substantially deeper in another direction
\cite{DoucotHohl2025,HienSabbah2015,Mochizuki2025}.  They determine transformed filtrations or
Stokes matrices.  The theorem here neither reconstructs those analytic data nor claims an analytic
Fourier equivalence; it isolates the formal valuation/contact invariant on the zero-loop,
possibly multilevel Airy locus.

\section{Leading-separated loopless cores}\label{sec:separated}

Let $I$ be an Airy tower of ramification $r$ and first divisor drop $r>e_1$.  Define
\begin{equation}\label{eq:first-factor}
 a=\frac r{e_1}\ge2.
\end{equation}
Its leading level is
\begin{equation}\label{eq:leading-slope}
 \frac{r+e_1}{r}=1+\frac1a.
\end{equation}

Two Stokes circles are \emph{leading-separated} if their common part is zero.  When their leading
slopes agree, this means that the leading coefficients lie in different Galois orbits; equivalently,
their appropriate powers differ.

\begin{theorem}[Separated-core formula]\label{thm:separated-core}
Let $I_i,I_j$ be leading-separated Airy towers with ramifications $r_i,r_j$ and first factors
$a_i,a_j$.  Then
\[
 B_{I_i,I_j}=\frac{r_ir_j}{\min(a_i,a_j)},
 \qquad
 \widetilde B_{I_i,I_j}=\frac1{\min(a_i,a_j)}.
\]
\end{theorem}

\begin{proof}
Suppose $a_i\le a_j$.  Then the leading slope of $I_i$ is at least that of $I_j$.  Since the
common part is zero, \eqref{eq:zero-common-edge} gives
\[
 B_{I_i,I_j}
 =r_j\bigl((r_i+e_{i,1})-r_i\bigr)
 =r_je_{i,1}
 =\frac{r_ir_j}{a_i}.
\]
The expression is symmetric when $a_i=a_j$, and the other ordering interchanges $i,j$.
\end{proof}

\begin{theorem}[Simultaneous realization]\label{thm:simultaneous}
Every finite family of strict divisor chains is realized by pairwise leading-separated Airy
towers.  Consequently every matrix supplied by Theorem~\ref{thm:separated-core} is a loopless
twisted core diagram.
\end{theorem}

\begin{proof}
For the $i$th chain choose
\[
 q_i=\sum_{h=1}^{m_i}c_{i,h}
 z^{(e_{i,h-1}+e_{i,h})/r_i},\qquad c_{i,h}\ne0.
\]
The gcd computation in Theorem~\ref{thm:factor-classification} shows that these are precisely the
prescribed level data.  For every fixed first factor $a$, there are only finitely many leading
coefficients to choose.  Choose their $a$th powers pairwise distinct.  Then equal-slope leading
terms lie in distinct $\mu_a$-orbits, while unequal slopes are automatically separated.  All
common parts are zero, and Theorems~\ref{thm:sign-law} and \ref{thm:separated-core} compute the
entire core.
\end{proof}

The reciprocal-min matrices automatically satisfy the required ultrametric condition.

\begin{proposition}[Separated ultrametric law]\label{prop:separated-ultra}
For positive integers $a,b,c$, the three values
\[
 \frac1{\min(a,b)},\qquad \frac1{\min(a,c)},\qquad \frac1{\min(b,c)}
\]
have their maximum at least twice.
\end{proposition}

\begin{proof}
The minimum of $a,b,c$ occurs in at least two of the three pairwise minima.  Taking positive
reciprocals reverses the order.
\end{proof}

The diagram forgets the factorization tail.

\begin{corollary}[Exact separated hidden fibre]\label{cor:hidden-fibre}
Fix labelled pairs $(r_i,a_i)$ in the leading-separated zero-loop locus.  The number of labelled
fission trees with these data and the same core matrix is
\[
 \prod_i A(r_i/a_i).
\]
The members have different level data and hence different admissible deformation classes.
\end{corollary}

\begin{proof}
After the first factor $a_i$ is fixed, the remaining divisor chain is exactly an ordered
factorization of $r_i/a_i$.  Theorem~\ref{thm:separated-core} ignores all these tails, while the
fission tree retains them.  Independence across labelled vertices gives the product.
\end{proof}

\section{Exact edge spectra on the elementary locus}\label{sec:spectrum}

\begin{definition}
An \emph{elementary loopless circle} is either the tame/unramified circle, recorded by $r=1$, or a
\emph{simple Airy circle} of ramification $r>1$, having the single level $(r+1)/r$.
\end{definition}

An elementary circle may contain additional unramified or lower inconsequential terms.  Such terms
do not change its sole level or its zero diagonal entry, but they control pairwise fission.

\begin{theorem}[Two-circle spectrum]\label{thm:pair-spectrum}
Let $I,J$ be distinct simple Airy circles of ramifications $r,s>1$.  Their possible edge
multiplicities are exactly
\[
 \Spec(r,s)=
 \begin{cases}
 \{\max(r,s)\}\cup rs\NN_{>0},&r\ne s,\\[2mm]
 \{0,1,\ldots,r\}\cup r^2\NN_{>0},&r=s.
 \end{cases}
\]
For mixed and unramified pairs,
\[
 \Spec(1,r)=\{1\}\cup r\NN_{>0}\quad(r>1),
 \qquad \Spec(1,1)=\NN_{\ge0}.
\]
Every listed value is realized.
\end{theorem}

\begin{proof}
All terms above a simple Airy level have integral exponent.  If the highest difference between
$I,J$ occurs at an integer $d\ge2$, the common terms above it have denominator one and make no
contribution to the rescaled edge formula.  The fission term gives
\[
 \widetilde B_{I,J}=d-1,
 \qquad B_{I,J}=rs(d-1),
\]
which realizes precisely the positive multiples of $rs$.

Assume no such high difference occurs.  If $r\ne s$, the slopes $1+1/r$ and $1+1/s$ differ.
The first difference is the larger slope and the common part is zero.  Formula
\eqref{eq:zero-common-edge} gives $B_{I,J}=\max(r,s)$.

Now let $r=s$.  If the mandatory leading coefficients lie in different $\mu_r$-orbits, the common
part is zero and \eqref{eq:zero-common-edge} gives $B_{I,J}=r$.  Otherwise make the leading terms
equal.  The first lower difference can occur exactly at
\[
 f=\frac jr,\qquad 1\le j\le r.
\]
Substitution in the common-level edge formula gives
\[
 B_{I,J}=rf-1=j-1,
\]
and hence $0,\ldots,r-1$.  Choosing distinct coefficients at the stated first differing exponent
realizes every value without changing the already complete ramification.  This proves the two
ramified formulas.

For an unramified and a ramified circle, the same high-difference argument gives $r\NN_{>0}$, while
separation at the Airy level gives $1$.  Two unramified circles satisfy
\eqref{eq:unramified-edge}, and a first difference at degree $d\ge1$ gives every $d-1\ge0$.
\end{proof}

\begin{corollary}[A local obstruction beyond ultrametricity]\label{cor:edge-three}
Two ramification-two simple Airy vertices have possible edge multiplicities
\[
 0,1,2,4,8,12,\ldots,
\]
but never $3$.  Since the ultrametric condition is vacuous on two vertices, this obstruction is
independent of the necessary condition \eqref{eq:known-ultra}.
\end{corollary}

\section{Finite fibres after ramification decoration}\label{sec:finite}

The preceding hidden fibres are finite.  This is part of a general finiteness statement.
We use the canonical full-data presentation of a labelled fission tree: fixed local
characteristics, finite common-prefix and sibling cases for every ordered pair, affine fission data
$(C,g,s_d)$, and the chosen finite compatibility predicate.  Thus the tree is, by definition, its
complete code satisfying those axioms; no quotient by unrecorded unary or coefficient-support data
is included in the fibre.

\begin{theorem}[Decorated finite-fibre theorem]\label{thm:finite-fibre}
Fix a finite labelled core matrix $B$ and a ramification $r_i\ge1$ at every vertex.  There are only
finitely many labelled fission trees whose decorated nonabelian Hodge diagram is $(B,r)$.  The
fibre has a complete finite enumeration.
\end{theorem}

\begin{proof}
At vertex $i$, the diagonal formula fixes
\begin{equation}\label{eq:fixed-weight}
 \sum_h\beta_{i,h}(e_{i,h-1}-e_{i,h})=B_{ii}+r_i^2-1=:W_i.
\end{equation}
If $W_i<0$, the fibre is empty.  Otherwise there are finitely many strict divisor chains from
$r_i$ to $1$.  Every summand on the left of \eqref{eq:fixed-weight} is positive, so every
$\beta_{i,h}$ lies in $[1,W_i]$.  Thus each vertex has finitely many possible level data.  Notice
that this step discards arbitrary coefficient support on inconsequential rays; fission trees record
level and fission data, not an individual polynomial representative.

Fix one level choice at every vertex.  For a pair of leaves, encode its common characteristic
prefix and sibling case by finite indices.  On each resulting interval write
$f_{ij}=s_d/r_i$.  The pair equation has the form
\[
  B_{ij}=C+s_dg-r_ir_j,
\]
where $C\ge0$ is fixed, $g\mid r_j$, and $g>0$.  Equivalently,
$B_{ij}+r_ir_j=C+s_dg$.  Positivity gives the explicit bounds
\[
 C\le B_{ij}+r_ir_j,\qquad g\le r_j,\qquad s_d\le B_{ij}+r_ir_j,
\]
and cancellation by $g>0$ shows that fixed $(C,g)$ determines at most one $s_d$.  Thus all raw
local and pair codes form a finite type.  A labelled fission tree is its complete raw code filtered
by the fission-tree compatibility predicate; every such subtype is finite.  Taking its universal
finite list gives the claimed complete enumeration.
\end{proof}

\begin{remark}
Ramification decoration is essential.  Without it, the one-vertex zero-loop graph is realized by a
simple Airy circle for every $r>1$, so its undecorated fibre is infinite.
\end{remark}

\section{Exact elementary decorated-graph recognition}\label{sec:recognition}

Extend the notation of Theorem~\ref{thm:pair-spectrum} by
\begin{equation}\label{eq:spec-all}
\begin{aligned}
 \Spec(1,1)&=\NN_{\ge0},\\
 \Spec(1,s)&=\{1\}\cup s\NN_{>0} &&(s>1),\\
 \Spec(r,s)&=\{\max(r,s)\}\cup rs\NN_{>0} &&(1<r\ne s>1),\\
 \Spec(r,r)&=\{0,1,\ldots,r\}\cup r^2\NN_{>0} &&(r>1).
\end{aligned}
\end{equation}

Let $N$ be a finite labelled set, let $B=(B_{ij})_{i,j\in N}$ be symmetric and nonnegative with
$B_{ii}=0$, and let $r_i\ge1$.  Put
\begin{equation}\label{eq:weights-recognition}
 w_{ij}=\frac{B_{ij}}{r_ir_j}\qquad(i\ne j),
 \qquad w_{ii}=0.
\end{equation}

\begin{theorem}[Elementary Stokes--Cartan recognition]\label{thm:recognition}
The following are equivalent.
\begin{enumerate}[label=(\roman*)]
\item $(B,r)$ is the decorated core diagram of a finite irregular class of elementary loopless
Stokes circles.
\item Both conditions hold:
\begin{align}
 B_{ij}&\in\Spec(r_i,r_j) &&(i\ne j),\label{eq:pair-test}\\
 w_{ij}\le w_{ik}\le w_{jk}&\Longrightarrow w_{ik}=w_{jk}
 &&(i,j,k\text{ pairwise distinct}).\label{eq:ultra-test}
\end{align}
\end{enumerate}
Moreover, the proof reconstructs explicit finite complex Puiseux polynomials simultaneously for
all labels.  The exhaustive unit-cost schedule consists of exactly $|N|^2+|N|^3$ ordered pair and
triangle queries, hence at most $2|N|^3$ queries when $N\ne\varnothing$.
\end{theorem}

Necessity follows immediately from Theorem~\ref{thm:pair-spectrum} and
\eqref{eq:known-ultra}.  We prove sufficiency in several steps.

\subsection{Finite ultrametric dendrograms}

We use the following standard form, included to fix the reconstruction.

\begin{lemma}[Dendrogram lemma]\label{lem:dendrogram}
Let $X$ be finite and let $d:X^2\to L$ be symmetric, zero on the diagonal, with values in a finite
totally ordered set.  Suppose no triangle has a unique largest side.  Then the values of $d$ are
realized as highest-difference heights in a rooted dendrogram.  At every height, distinct child
clusters can be assigned distinct coefficients in any infinite field.
\end{lemma}

\begin{proof}
Let $D$ be the largest value.  Define $x\sim y$ when $x=y$ or $d(x,y)<D$.  If
$d(x,y),d(y,z)<D$ but $d(x,z)=D$, the triangle has a unique largest side.  Thus $\sim$ is an
equivalence relation.  Attach distinct coefficients to its finitely many classes at height $D$ and
recurse inside each class.  The number of distinct values decreases, so the construction
terminates.  Two leaves first differ at the coefficient attached to their least common ancestor,
whose height is their prescribed value.
\end{proof}

Equivalently, for every threshold $t\ge0$, the relation $d(x,y)\le t$ is an equivalence relation,
and these partitions coarsen with $t$.  We use this threshold formulation below.

\subsection{Integral heights}

Assume \eqref{eq:pair-test}--\eqref{eq:ultra-test}.  Apply
Lemma~\ref{lem:dendrogram} to the positive integral values among the $w_{ij}$, from largest to
smallest.  By \eqref{eq:spec-all}, every such value $n$ arises from the high-contact branch of the
pair spectrum and is realized by a first difference at the unramified exponent $n+1$.  At the
corresponding dendrogram node assign distinct complex coefficients to its child clusters at
$z^{n+1}$.  Integer terms are admissible and inconsequential for every ramification.  After all
such nodes are processed, every unresolved pair has weight $<1$.

\subsection{Mandatory Airy levels}

Consider one unresolved cluster.  Let $r>1$ be the smallest ramification among its ramified leaves;
regard unramified leaves as having ramification $\infty$.  Pair-spectrum membership forces:
\begin{enumerate}[label=(\roman*)]
\item an $r$-leaf and any leaf of larger ramification have weight exactly $1/r$;
\item two $r$-leaves have weight at most $1/r$;
\item two leaves of larger ramification have weight strictly less than $1/r$.
\end{enumerate}
Indeed, after integral weights have been removed, only the low branches of
\eqref{eq:spec-all} remain.

It follows from \eqref{eq:ultra-test} that the equivalence classes for $w<1/r$ consist of some
blocks of $r$-leaves and at most one block containing all larger-ramification and unramified leaves.
Assign to the $r$-blocks nonzero coefficients at
\[
 z^{1+1/r}
\]
whose $r$th powers are pairwise distinct.  Assign coefficient zero to the possible residual block.
Thus mandatory siblings are separated modulo the $\mu_r$ Galois action, and there is exactly the
single possible empty child allowed in the twisted fission-tree sibling axioms
\cite[\S3.4]{BoalchDoucotRembado2025}.  Remove the $r$-blocks and repeat with the next ramification
inside the residual block.

\subsection{Lower fractional heights}

Inside one fixed $r$-block, distinct leaves have equal mandatory Airy term.  Their unresolved
weights are
\[
  \frac b{r^2},\qquad 0\le b\le r-1.
\]
At the largest numerator $b$, the dendrogram lemma partitions by smaller weight.  Assign distinct
coefficients to the child clusters at the exponent
\[
  \frac{b+1}{r}.
\]
The computation in Theorem~\ref{thm:pair-spectrum} gives edge multiplicity $b$ between distinct
children.  Recurse on $b$.  At $b=0$, distinct coefficients of $z^{1/r}$ realize edge zero.  These
terms introduce no new level: the mandatory exponent $(r+1)/r$ has already reduced the cumulative
gcd to one.

After all ramified blocks are resolved, any residual unramified block has pair weight zero.
Distinct coefficients of $z$ realize it.

\begin{proof}[Completion of Theorem~\ref{thm:recognition}]
The preceding finite recursion constructs Puiseux polynomials $q_i$.  When $r_i>1$, the only gcd
drop is the mandatory exponent $(r_i+1)/r_i$, so $\langle q_i\rangle$ is a simple Airy circle and
has zero diagonal by Theorem~\ref{thm:sign-law}.  When $r_i=1$, it is unramified and loopless.

The construction may be written without recursive coefficient choices.  Use the finite global
stage set consisting of high stages $k+1$, mandatory stages $1+1/r$, and lower stages
$(b+1)/r$.  At each stage, assign to every strict-threshold equivalence class a distinct positive
integer complex coefficient.  The acute-isosceles law makes these classes transitive.  The exact
pair spectrum selects, for each of its nine oriented mechanisms, a highest separating stage.

Equal rational exponents from different stages could in principle cancel after collection.  They
do not: high, mandatory and lower exponent ranges are disjoint except for inactive stages, while
within one range the ramification and numerator identities force equality of the stages.  Thus the
selected stage is exponent-isolated.  Summing all stages gives actual finite Puiseux polynomials
$q_i$ whose coefficient at the selected exponent differs and whose coefficients at every larger
exponent agree.  Hence every pair realizes exactly $B_{ij}$ simultaneously.  Certificate-form
sufficiency follows.  The exhaustive query set consists of all ordered pairs and ordered triples,
so it has exactly $|N|^2+|N|^3$ elements and at most $2|N|^3$ elements when $N\ne\varnothing$.
\end{proof}

\begin{corollary}\label{cor:no-further-obstruction}
On the elementary locus, the local spectra \eqref{eq:spec-all} and rescaled ultrametricity are a
complete set of obstructions.  In particular, there are infinitely many forbidden two-vertex and
higher decorated graphs not detected by ultrametricity alone.
\end{corollary}

\begin{proof}
Completeness is Theorem~\ref{thm:recognition}.  For the infinite two-vertex family, two
ramification-two vertices with edge $4k+3$ fail the exact spectrum for every $k\ge0$.  Triangle
conditions are vacuous on two labels.
\end{proof}

\section{Examples and information loss}\label{sec:examples}

\begin{example}[An all-ramified graph twin]\label{ex:twin}
Consider the Airy towers
\[
\begin{array}{ccl}
 r=2:&2>1,&q_1=c_1z^{3/2},\\
 r=3:&3>1,&q_2=c_2z^{4/3},\\
 r=4:&4>2>1,&q_3=c_3z^{3/2}+d_3z^{3/4}.
\end{array}
\]
Choose $c_1,c_3$ in different $\mu_2$-orbits.  The first factors are $2,3,2$, so
Theorem~\ref{thm:separated-core} gives
\[
 B_{12}=3,\qquad B_{13}=4,\qquad B_{23}=6.
\]
This asserts only the displayed decorated core data, not an isomorphism of associated analytic
moduli spaces.
\end{example}

\begin{example}[Exponential hidden fibres]
Fix a separated vertex of ramification $p^k$ and first factor $p^j$.  The graph and first
ramification drop forget the remaining factorization of $p^{k-j}$, leaving
$2^{k-j-1}$ possibilities when $k-j\ge1$ and one possibility when $k=j$.  Products of such
vertices give exponentially large labelled fibres with an unchanged core matrix.  Indeed,
Corollary~\ref{cor:counts} counts the remaining ordered factorizations, while
Theorem~\ref{thm:separated-core} shows that later factors do not change the separated core.
\end{example}

The finite-fibre theorem and these examples isolate the role of decoration.  The undecorated Cartan
shadow can hide arbitrarily large ramification and infinitely many local types, as the simple Airy
family already shows.  Ramification makes
the full fibre finite; on the separated locus, the first quotient fixes the matrix and the remaining
ordered factors describe the exact hidden part.  This complements existing counts in the untwisted
case \cite{Boalch2025Counting}: the present counts concern one ramified zero-loop stratum rather
than untwisted trees of fixed slope and leaf number.

\section{Scope}\label{sec:scope}

Several stronger statements are intentionally excluded.
\begin{enumerate}[label=(\roman*)]
\item We do not classify all twisted loopless vertices or all twisted nonabelian Hodge graphs.
\item Equal diagrams do not imply isomorphic wild character varieties, isomonodromy systems, or
nonabelian Hodge spaces.  Enriched-tree data are needed for such comparisons
\cite{Doucot2026Basic}.
\item Knot theory appears only as contextual comparison; no knot-theoretic assertion enters a
paper theorem.
\item The exact $n^2+n^3\le2n^3$ recognition count is a unit-cost query bound, not a bit-complexity
analysis of rational arithmetic.
\end{enumerate}

\section{Concluding perspective}

The universal law is formal-contact covariance: tangent changes of uniformizer preserve valuation,
and therefore preserve every invariant assembled solely from leading orders and pairwise
differences.  Sheet coarea, complementary contact, threshold tomography, additive energy and the
negative Fourier rank chart are shadows of this one law.  The remaining arithmetic theorems are
its zero-loop collateral: imposing the critical diagonal value discretizes the valuation data into
the gcd gap, Airy divisor chains and exact pair spectra.  Thus coordinate covariance explains why
the geometric skeleton survives Fourier motion, while arithmetic rigidity explains which integral
Cartan shadows can occur.

The equality $B_{I,I}=0$ looks at first like one linear constraint on a Puiseux characteristic.
The gcd lattice makes it rigid: each characteristic exponent is forced to be the sum of two
successive ramifications.  This converts a Stokes-geometric equality case into divisor chains,
ordered factorizations, negative Fourier steps and iterated cabling data.

At graph level, the same rigidity has two opposite effects.  It creates large hidden fibres because
the Cartan shadow sees only the first factor of a separated Airy tower.  Yet it also creates sharp
edge gaps: local divisibility obstructs many graphs that pass every ultrametric test.  On the
elementary locus the two effects fit exactly.  Local spectra supply the missing arithmetic
conditions, and the ultrametric reconstructs their global fission dendrogram.  This gives a first
complete recognition theorem for a genuinely ramified, arbitrary-size class of nonabelian Hodge
graphs.

\begin{thebibliography}{99}

\bibitem{BiquardBoalch2004}
O.~Biquard and P.~Boalch,
\emph{Wild non-abelian Hodge theory on curves},
Compos. Math. \textbf{140} (2004), no.~1, 179--204.
\href{https://doi.org/10.1112/S0010437X03000010}{doi:10.1112/S0010437X03000010}.

\bibitem{Boalch2025Counting}
P.~Boalch,
\emph{Counting the fission trees and nonabelian Hodge graphs (untwisted case)},
J. Geom. Phys. \textbf{214} (2025), 105512.
\href{https://doi.org/10.1016/j.geomphys.2025.105512}{doi:10.1016/j.geomphys.2025.105512}.

\bibitem{BoalchDoucotRembado2025}
P.~Boalch, J.~Dou\c cot and G.~Rembado,
\emph{Twisted local wild mapping class groups: configuration spaces, fission trees and complex
braids},
Publ. Res. Inst. Math. Sci. \textbf{61} (2025), no.~3, 391--445.
\href{https://doi.org/10.4171/PRIMS/61-3-4}{doi:10.4171/PRIMS/61-3-4}.

\bibitem{BoalchYamakawa2020}
P.~Boalch and D.~Yamakawa,
\emph{Diagrams for nonabelian Hodge spaces on the affine line},
C. R. Math. \textbf{358} (2020), no.~1, 59--65.
\href{https://doi.org/10.5802/crmath.11}{doi:10.5802/crmath.11}.

\bibitem{Doucot2025Graphs}
J.~Dou\c cot,
\emph{New nonabelian Hodge graphs from twisted irregular connections},
arXiv:2509.24861 (2025), revised 2026.
\href{https://arxiv.org/abs/2509.24861}{arXiv:2509.24861}.

\bibitem{Doucot2026Basic}
J.~Dou\c cot,
\emph{Basic representations of genus zero nonabelian Hodge spaces},
SIGMA \textbf{22} (2026), 027, 43 pp.
\href{https://doi.org/10.3842/SIGMA.2026.027}{doi:10.3842/SIGMA.2026.027}.

\bibitem{Doucot2026Fourier}
J.~Dou\c cot,
\emph{Fourier transform of irregular connections on $\mathbb P^1$ and classification of
Argyres--Douglas theories},
arXiv:2603.15942 (2026).
\href{https://arxiv.org/abs/2603.15942}{arXiv:2603.15942}.

\bibitem{DoucotHohl2025}
J.~Dou\c cot and A.~Hohl,
\emph{A topological algorithm for the Fourier transform of Stokes data at infinity},
J. Lond. Math. Soc. \textbf{112} (2025), e70253.
\href{https://doi.org/10.1112/jlms.70253}{doi:10.1112/jlms.70253}.

\bibitem{GarciaBarrosoGonzalezPopescu2019}
E.~R.~Garc\'ia Barroso, P.~D.~Gonz\'alez P\'erez and P.~Popescu-Pampu,
\emph{The valuative tree is the projective limit of Eggers--Wall trees},
RACSAM \textbf{113} (2019), no.~4, 4051--4105.
\href{https://doi.org/10.1007/s13398-019-00646-z}{doi:10.1007/s13398-019-00646-z}.

\bibitem{Hiroe2017}
K.~Hiroe,
\emph{Ramified irregular singularities of meromorphic connections and plane curve singularities},
Josai Math. Monogr. \textbf{10} (2017), 161--192.
\href{https://arxiv.org/abs/1406.5788}{arXiv:1406.5788}.

\bibitem{HienSabbah2015}
M.~Hien and C.~Sabbah,
\emph{The local Laplace transform of an elementary irregular meromorphic connection},
Rend. Semin. Mat. Univ. Padova \textbf{134} (2015), 133--196.
\href{https://doi.org/10.4171/RSMUP/134-4}{doi:10.4171/RSMUP/134-4}.

\bibitem{Mochizuki2025}
T.~Mochizuki,
\emph{Stokes shells and Fourier transforms},
arXiv:1808.01037, version 4 (2025).
\href{https://arxiv.org/abs/1808.01037}{arXiv:1808.01037}.

\enlargethispage{3\baselineskip}
\end{thebibliography}

\end{document}
